% Full answers PDF: MCQ answers + worked explanations per option, then official MS clips.
% Compiler: XeLaTeX.
\documentclass[12pt,a4paper]{article}
\usepackage[margin=18mm,headheight=14pt]{geometry}
\usepackage{fontspec}
\usepackage[version=4]{mhchem}
\usepackage{graphicx}
\usepackage{xcolor}
\usepackage{booktabs}
\usepackage{fancyhdr}
\usepackage{enumitem}
\usepackage{needspace}

\setmainfont{Helvetica Neue}

\pagestyle{fancy}
\fancyhf{}
\fancyhead[L]{\small 3.3 Metallic bonding --- answers}
\fancyhead[R]{\small CIE 9701}
\fancyfoot[C]{\thepage}
\renewcommand{\headrulewidth}{0.3pt}

\setlength{\parindent}{0pt}
\setlength{\parskip}{0.35em}

\newcommand{\msblock}[2]{%
  \par\vspace{0.8em}%
  \noindent\begin{minipage}{\textwidth}%
  {\normalsize\bfseries #1}\par\vspace{0.25em}%
  \noindent\includegraphics[width=\textwidth,keepaspectratio]{#2}\par
  \end{minipage}\par
}

\newcommand{\mcq}[4]{%
  \par\needspace{5\baselineskip}\vspace{0.7em}%
  {\normalsize\bfseries #1.\quad #2}%
  \hfill{\footnotesize\color{gray}#3}\par\vspace{0.15em}%
  {\small #4}\par
}

\newcommand{\why}[2]{\par\needspace{2.2\baselineskip}\vspace{0.1em}{\small\textbf{#1} #2}\par}

\begin{document}

\begin{center}
{\large\bfseries 3.3 Metallic bonding homework --- answers}\\[0.2em]
{\small CIE 9701 AS Chemistry \quad·\quad 8 multiple-choice + 11 written marks}
\end{center}

{\small\bfseries Do not issue this file to students.}
\hrule
\vspace{0.4em}

\section*{Section A --- Multiple choice}

{\normalsize\bfseries Quick answers}

\vspace{0.3em}
{\small
\begin{tabular}{@{}ll@{\hspace{2.6em}}ll@{\hspace{2.6em}}ll@{}}
\toprule
Q & Ans. & Q & Ans. & Q & Ans.\\
\midrule
A1 & \textbf{A} & A4 & \textbf{A} & A7 & \textbf{B}\\
A2 & \textbf{A} & A5 & \textbf{A} & A8 & \textbf{D}\\
A3 & \textbf{A} & A6 & \textbf{A} & & \\
\bottomrule
\end{tabular}}

\vspace{0.6em}
{\normalsize\bfseries Worked explanations}

\mcq{A1}{A}{2023 MJ Paper 11 Q2}{%
Both calcium and barium are Group 2 metals with the same metallic structure, so
the only difference that matters is the size of the cation. \ce{Ca^2+} is smaller
than \ce{Ba^2+}, so the delocalised electrons are held closer to the \ce{Ca^2+}
ions and the electrostatic attraction is stronger, giving calcium the higher
melting point.}

\why{B:}{metals do not have a ``partly giant molecular'' structure --- they have a
giant \emph{metallic} lattice of cations in a sea of delocalised electrons.}

\why{C:}{both Ca and Ba contribute the \emph{same} number (two) of delocalised
electrons per atom; ionisation energy is not the factor here.}

\why{D:}{there is no such repulsion --- the attraction between the cations and
the delocalised electrons is what holds the lattice together.}

\mcq{A2}{A}{2021 MJ Paper 13 Q36}{%
All three statements are correct and together explain the rise in melting point
from Na to Al. Going Na $\rightarrow$ Mg $\rightarrow$ Al:
(1) the charge on the metal ion increases (\ce{Na+}, \ce{Mg^2+}, \ce{Al^3+});
(2) the number of delocalised electrons per ion increases (1, 2, 3);
(3) the ionic radius decreases.
A higher charge, more delocalised electrons and a smaller radius all strengthen
the metallic bonding, so the melting point rises.}

\why{B / C:}{each drops one of the three correct statements.}

\why{D:}{keeps only statement 1 --- the other two statements are also true.}

\mcq{A3}{A}{2024 MJ Paper 12 Q6}{%
Boron in \ce{BF3} has only \textbf{three} bond pairs and \textbf{no lone pair},
so there are only 6 electrons in its outer shell --- an \emph{incomplete} octet.
This is why \ce{BF3} is a good electron-pair acceptor (a Lewis acid).}

\why{B:}{in \ce{CH3-} the carbon has 3 bond pairs \emph{and} one lone pair =
8 outer electrons (complete octet).}

\why{C:}{in \ce{F2O} the oxygen has 2 bond pairs and 2 lone pairs = 8 electrons.}

\why{D:}{in \ce{H3O+} the oxygen has 3 bond pairs and 1 lone pair = 8 electrons.}

\mcq{A4}{A}{2021 FM Paper 12 Q34}{%
All three species have a vacant orbital that can accept a lone pair, so all three
can form a coordinate (dative covalent) bond.
\ce{BF3}: B has an empty 2p orbital.
\ce{H+}: the hydrogen ion has an empty 1s orbital.
\ce{CH3+}: the carbon has an empty 2p orbital (only 6 outer electrons).}

\why{B / C / D:}{each option excludes at least one species that is in fact able
to accept a lone pair. The trap is forgetting that \ce{H+} and \ce{CH3+}, not
just \ce{BF3}, are electron-deficient.}

\mcq{A5}{A}{2024 MJ Paper 11 Q5}{%
In the \ce{Al2Cl6} dimer the two aluminium atoms are joined through two
\emph{bridging} chlorine atoms. Each bridging Cl uses a lone pair to form a
\emph{coordinate (dative covalent)} bond to one Al, and the Al--Cl bonds are
covalent. So \ce{Al2Cl6} contains \textbf{covalent and coordinate} bonding only.}

\why{B:}{``covalent only'' ignores the dative bonds to the bridging Cl atoms ---
each Al is surrounded by 4 electron pairs, which needs the dative bonds.}

\why{C / D:}{\ce{Al2Cl6} is a molecular (covalent) species; it is not ionic. (It
is \ce{AlCl3} that behaves ionically in the solid state, not the \ce{Al2Cl6}
dimer.)}

\mcq{A6}{A}{2023 ON Paper 11 Q5}{%
Aluminium chloride \emph{sublimes} to give a vapour of \textbf{\ce{Al2Cl6}
dimer} molecules --- two Al centres bridged by two Cl atoms, with four terminal
Cl atoms. Option A is the dimer.}

\why{B / C:}{do not show the correct \ce{Al2Cl6} dimer arrangement (the vapour is
not the simple \ce{AlCl3} monomer at this temperature).}

\why{D:}{shows an ionic formulation, \ce{Al^3+(Cl-)3}. The vapour consists of
covalent \ce{Al2Cl6} molecules, not separate ions.}

\mcq{A7}{B}{2022 FM Paper 12 Q5}{%
The correct dot-and-cross diagram must show each aluminium atom surrounded by
\textbf{four} electron pairs (giving Al a complete octet) and must show the
\emph{dative} bond: a bridging chlorine atom donates a lone pair to an aluminium
atom. The key on the paper distinguishes the electrons that came from Al from
those that came from Cl. Diagram B is the one that shows the \ce{Al2Cl6}
structure with the bridging Cl atoms donating a lone pair, and the correct
number of electron pairs around each Al.}

\why{A:}{shows \ce{Al^3+} and \ce{Cl-} ions with the wrong electron bookkeeping
--- \ce{Al2Cl6} is molecular and each Al still has four shared pairs.}

\why{C / D:}{show the wrong arrangement of atoms / electron pairs (they do not
reproduce the bridged \ce{Al2Cl6} structure with the dative pairs correctly
marked).}

\mcq{A8}{D}{2021 ON Paper 11 Q35}{%
A chloride with \emph{both} covalent and coordinate bonding must be a molecular
chloride that forms a bridged dimer.
(1) \textbf{Al} --- forms \ce{Al2Cl6}, which has covalent bonds and dative bonds
from the bridging Cl atoms. \textbf{Correct.}
(2) Si --- forms \ce{SiCl4}, which has ordinary covalent bonds only (no dative
bonding). \textbf{Wrong.}
(3) Mg --- forms \ce{MgCl2}, which is ionic. \textbf{Wrong.}
Only statement 1 is correct.}

\why{A / B / C:}{all include Si or Mg, whose chlorides do not contain coordinate
bonding.}

\newpage

\section*{Section B --- Mark scheme clips}

\msblock{B1(a)  [2]}{cie-9701-2025-fm-p22-q2a-ms.png}

\msblock{B2(c)  [1]}{cie-9701-2024-mj-p22-q1c-ms.png}

\msblock{B3(a)  [1]}{cie-9701-2024-mj-p21-q5a-ms.png}

\msblock{B4(b)(i)  [1]}{cie-9701-2021-mj-p22-q2b-i-ms.png}

\msblock{B4(b)(ii)  [2]}{cie-9701-2021-mj-p22-q2b-ii-ms.png}

\msblock{B5(a)  [2]}{cie-9701-2024-mj-p23-q4a-ms.png}

\msblock{B6(e)(i)  [2]}{cie-9701-2021-fm-p22-q2e-i-ms.png}

\end{document}
